%% ch06 --- The butterfly effect, measured
%%
%% Stub, generated by tools/gen_stubs.py from tools/chapters.json.
%% The brief below is the contract for this chapter: what it must argue, what
%% it must buy the reader, and what it must NOT take from its neighbours.
%% Writing the chapter means deleting the stub block below (the one that opens
%% on the next \chapter line) and replacing it with the chapter text -- after
%% which this file is never regenerated.
%% If the brief turns out to be wrong once the code has been run, change it in
%% the manifest and record the contradiction in CLAUDE.md under *Resolved
%% questions*.

\chapter{The butterfly effect, measured}
\label{ch:butterfly}

\chapterstub{%
Argument: sensitive dependence on initial conditions is not a metaphor but a
measurement. Two runs of the logistic map at r = 4 that start one part in
ten billion apart stay indistinguishable for a while, then separate
completely within a few dozen steps \dash{} counted exactly, in pure Python
arithmetic, so the number of steps on the page is the same on every machine.
The separation roughly doubles each step, so the error grows exponentially,
which is Chapter 2's growth law now working against us. Edward Lorenz's 1961
discovery, when he restarted a run from a printout rounded to three decimals
(0.506 for 0.506127) and got a different weather, and the title of his 1972
talk about a butterfly in Brazil. Determinism and predictability are two
different properties: the map is perfectly deterministic \dash{} the same
start gives the same future, checked in code \dash{} and still unpredictable
in practice, because no start is known exactly. Payoff: the reader can
measure how fast two nearby starts separate and say, for a given measurement
error, how many steps a prediction is good for. Traps to elicit: chaos means
random; the butterfly causes the tornado; if I rerun it I will get something
different. Listings: two nearby runs side by side; the separation table;
Lorenz's rounding reproduced on the logistic map. Labs: find the step at
which two runs first differ by more than a tolerance. Figures: two
trajectories peeling apart; separation on a log scale, a straight line then
a ceiling.

\medskip
\textbf{Word budget:} 3000.
}
